\begin{frame}{반복하면 미는 점이 3개만 남는다}
\vfill
\centering\includegraphics[height=0.74\textheight]{sv_train4.png}\par\vskip0.2em
{\large 1걸음 뒤로 미는 점은 처음부터 \kb{서포트 벡터 3개}뿐}
\vfill
\note{2만 걸음 중 몇 장면입니다. 금색 동그라미가 그 걸음에 가장 가까워 미는 점이고, 길 가장자리는 경계에서 m만큼 떨어진 점선입니다. 1걸음 뒤 첫 선에서는 빨간 점 0.7, 0.2이 가장 가깝습니다. 그 뒤로 빨강 0.7, 0.2과 파랑 마이너스 0.6, 마이너스 0.4이 번갈아 가장 가까워지며 m이 0.62, 0.66, 0.64, 0.67, 0.65, 0.67로 조금씩 흔들립니다. 9걸음 뒤에는 파란 점 0.1, 마이너스 1.1이 처음으로 가장 가까워집니다. 2만 걸음 뒤에는 이 서포트 벡터 3개가 번갈아 밀며 m 0.67의 길을 붙잡고 있습니다. 1걸음 뒤로 한 번이라도 민 점은 이 3개뿐이고, 나머지 17개는 한 번도 밀지 않았습니다.}
\end{frame}